\documentclass[11pt,a4paper]{article}
\usepackage[margin=1in]{geometry}
\usepackage{amsmath,amssymb}
\usepackage{booktabs}
\usepackage{array}
\usepackage{listings}
\usepackage{xcolor}
\definecolor{seclink}{HTML}{1F4E8C}
\usepackage[colorlinks=true,linkcolor=seclink,citecolor=seclink,urlcolor=blue]{hyperref}
\usepackage{parskip}
\usepackage[most]{tcolorbox}
\usepackage{tikz}
\usetikzlibrary{positioning,arrows.meta,fit,backgrounds,calc}

\definecolor{ink}{HTML}{1A1A1A}
\definecolor{grey}{HTML}{5F5B54}
\definecolor{pale}{HTML}{F5F3EE}
\definecolor{yellowbg}{HTML}{FBF3D5}
\definecolor{pinkbg}{HTML}{FADDE1}
\definecolor{tensor}{HTML}{1F6F4A}
\definecolor{seq}{HTML}{1F4E8C}
\definecolor{prod}{HTML}{9C3B1B}
\definecolor{sumc}{HTML}{6B3FA0}
\definecolor{codebg}{RGB}{246,246,244}

\tcbset{plainbg/.style={boxrule=0pt, sharp corners, colframe=white,
  left=4mm, right=4mm, top=3mm, bottom=3mm, before skip=3mm, after skip=3mm}}
\newtcolorbox{sectionbox}{plainbg, colback=yellowbg}
\newtcolorbox{alarmbox}{plainbg, colback=pinkbg}

\tikzset{
  lvl/.style={draw=ink, fill=white, rounded corners=2.5pt, align=center,
              font=\small\ttfamily, inner sep=4.5pt, minimum width=44mm,
              minimum height=8mm},
  leafn/.style={draw=ink, fill=white, rounded corners=2pt, align=center,
                font=\scriptsize\ttfamily, inner sep=3.4pt, minimum height=6mm},
  op/.style={circle, draw, inner sep=0.6pt, minimum size=7mm, font=\small\bfseries},
  ar/.style={-{Stealth[length=1.8mm]}, semithick, color=grey},
  lbl/.style={font=\scriptsize\itshape, color=grey},
}

\lstdefinestyle{code}{
  basicstyle=\ttfamily\small,
  commentstyle=\color{grey}\itshape,
  breaklines=true,
  frame=leftline,
  framesep=7pt,
  xleftmargin=12pt,
  rulecolor=\color{gray!45},
  backgroundcolor=\color{codebg},
  columns=fullflexible,
  showstringspaces=false,
  keepspaces=true,
  literate=
    {↓}{{$\downarrow$}}1 {↑}{{$\uparrow$}}1
    {→}{{$\rightarrow$}}1 {⊠}{{$\boxtimes$}}1
    {◁}{{$\lhd$}}1 {×}{{$\times$}}1 {—}{{\textemdash}}1 {’}{{'}}1,
}
\lstset{style=code}

\newcommand{\seqb}{\lhd^{\ast}}


\title{\vspace{-2em}Interestingness as an Objective}
\author{\normalsize Research notes on reward, scarcity, and machine-tuned game design}
\date{September 2026}

\begin{document}
\maketitle
\vspace{-2em}

\begin{sectionbox}
The Labyrinth is a turn-based dungeon crawler whose levels are generated and whose monsters are
language models. An early draft of it had no strategy and no pleasure in it, and the fault was
not in the writing. It was that the game had operators and nothing scarce behind them.

\medskip

This report is the evidence for the repair. It asks two questions. What is actually known about
why a resource-limited turn-based game feels good, and can a machine be given something to
maximise that stands in for feeling good. The answer to the second is yes, with a documented
existence proof from 2011, and it is not difficulty and it is not one number.

\medskip

The central claim arrived last and reorganises the rest. Two measurements are needed and neither
implies the other. \emph{Depth}, meaning bad play is punished, and \emph{breadth}, meaning that
among good play there is no best. A game with depth and no breadth is a puzzle. The accompanying
game \texttt{stalin} has both, by holding three ends apart and refusing to say which is the real
one, and it is the working model for everything in Section~\ref{sec:plural}.
\end{sectionbox}

\section{The two questions}

They are usually confused and they are not the same.

The first is a design question. Given a turn-based game with limited resources and one player,
which structures produce the experience of playing well rather than the experience of being
processed. That question has a literature in psychology, a much larger literature in the practice
of the genre, and a small number of findings sharp enough to act on.

The second is an engineering question. The design of the game leaves ten or so numbers
undetermined, and choosing them by argument is choosing them badly. Can a program play the game
many times and choose them by evidence. That question has a literature in automated game design,
a published success, and a cost problem that turns out to be solvable here for a reason peculiar
to this architecture.

\begin{sectionbox}
The connection between them is the reason for one report rather than two. A tuner needs an
objective function, and an objective function is a design claim written as arithmetic. So the
first question is not preliminary to the second. It \emph{is} the second, done carefully enough
to compile.
\end{sectionbox}

\section{Reward, stated correctly}
\label{sec:reward}

The word dopamine is used loosely in discussions of games and it is worth being precise, because
the precise version says something useful and the loose version says only that pleasure exists.

\subsection{The signal is an error, not a reward}

Midbrain dopamine neurons do not fire in proportion to reward. They fire in proportion to
\emph{reward prediction error}, which is the difference between what arrived and what was
expected\footnote{Schultz, \emph{Dopamine reward prediction error coding}, Dialogues in Clinical
Neuroscience 18(1), 2016. See also Schultz, \emph{Dopamine reward prediction-error signalling: a
two-component response}, Nature Reviews Neuroscience 17, 2016.}. A fully predicted reward
produces almost no phasic response. An unexpected one produces a large positive response, and an
expected one that fails to arrive produces a dip below baseline.

Three consequences follow directly and all three are design constraints.

\begin{center}
\begin{tabular}{@{}l@{\hspace{5mm}}>{\raggedright\arraybackslash}p{0.62\textwidth}@{}}
\toprule
\textbf{a prediction is required} & A reward the player did not see coming cannot be better than
expected, because nothing was expected. The player must have been estimating an outcome, which
means the game must have given them the means to estimate it. \\[1mm]
\textbf{predictable is silent}    & Loot drawn from a fixed table by depth is fully predicted and
therefore produces nothing. This is what the first draft of the design did. \\[1mm]
\textbf{the dip is real}          & An expected reward that fails to arrive is negatively
signalled. Loss has a mechanism, and it is not merely the absence of gain. \\
\bottomrule
\end{tabular}
\end{center}

The first of those is the one usually missed. Surprise is not the same as randomness. A player
who cannot form a prediction at all is not surprised by the outcome, they are indifferent to it,
and the way to make outcomes surprising is to make the situation legible enough that the player
commits to a guess.

\subsection{The finding worth building on}

There is a second component to the dopamine response, discovered by Fiorillo, Tobler and
Schultz, and it is the single most useful quantitative result in this
report\footnote{Fiorillo, Tobler and Schultz, \emph{Discrete Coding of Reward Probability and
Uncertainty by Dopamine Neurons}, Science 299, 2003.}.

Alongside the error signal there is a gradual ramp in activity between the cue and the moment the
reward is due, and its size tracks \emph{uncertainty} rather than value. It is maximal when the
probability of reward is one half, and it falls away as the probability approaches either zero or
one.

\begin{sectionbox}
Which gives a target rather than a direction. The advice to make a game uncertain is not
actionable, because uncertainty has no natural units. The advice to arrange that the pivotal
decision of a run sits near even odds is actionable, is measurable from a batch of playthroughs,
and is a legitimate objective for an optimiser.
\end{sectionbox}

Two honest qualifications. The measurements are of single neurons in macaques anticipating juice,
and the distance from there to a player deciding whether to descend one more level is very
large. And a design rule near one half is independently supported by ordinary design reasoning
and by Browne's uncertainty criterion in Section~\ref{sec:browne}, which arrived at it from
self-play statistics with no neuroscience involved. So the correct reading is convergence
between three arguments rather than derivation from one, and the rule is worth adopting for that
reason and not because it has a brain scan attached.

\subsection{What this does not license}
\label{sec:notlicensed}

The literature on gamification is largely a record of this material being misused, and the
failure mode is well documented enough to state as a warning.

The overjustification effect is that supplying an expected external reward for an activity which
was already intrinsically interesting reduces the motivation to do it. Reviews of applied
gamification report that arbitrary extrinsic rewards do not merely fail to help, they measurably
displace the motivation that was there\footnote{Dah et al., \emph{Gamification is not Working:
Why?}, Games and Culture, 2025. See also the overjustification literature descending from Deci
and Ryan.}. Points, badges and leaderboards bolted onto something already enjoyable make it less
enjoyable once the points stop meaning anything.

\begin{alarmbox}
So the instruction here is narrow. Nothing in this report recommends adding a reward layer to the
Labyrinth. Every recommendation is about making the decisions the player already makes have
consequences they can estimate. The pleasure is meant to come from the estimate turning out to
be right, which is intrinsic to playing, and not from a number going up beside the screen.
\end{alarmbox}

\subsection{Autonomy and competence}

The complementary positive account is self-determination theory, applied to games by Ryan, Rigby
and Przybylski across four studies\footnote{Ryan, Rigby and Przybylski, \emph{The Motivational
Pull of Video Games: A Self-Determination Theory Approach}, Motivation and Emotion 30, 2006. See
also Przybylski, Rigby and Ryan, \emph{A Motivational Model of Video Game Engagement}, Review of
General Psychology 14, 2010.}. They find that perceived autonomy and perceived competence
independently predict enjoyment, preference and post-play well-being, and that they do so
independently of the game's content.

Read as design constraints, autonomy means the game must offer choices whose alternatives were
genuinely available, and competence means the player must be able to tell that the outcome
followed from their choice. Both are properties of the decision structure and neither is a
property of the reward schedule.

\begin{sectionbox}
Which is convenient, because both are exactly what the operator census of the Labyrinth
measures. A tensor offers no autonomy, because both sides happen regardless. A product with
nothing scarce behind it offers no competence, because the outcome did not follow from the
choice. The two psychological requirements have algebraic signatures.
\end{sectionbox}

\section{Intermittent reward}
\label{sec:intermittent}

The mechanism is operant conditioning and the relevant schedule has a name. Under a
\emph{variable ratio} schedule the reward arrives after an unpredictable number of responses, and
that schedule is the most productive and the most resistant to extinction of any that
Skinner classified\footnote{Standard treatment in Simply Psychology, \emph{Schedules of
Reinforcement}, and Explore Psychology, \emph{Variable Ratio Schedule of Reinforcement}. The
original taxonomy is Ferster and Skinner, \emph{Schedules of Reinforcement}, 1957.}. Resistance to
extinction is the technical name for what is meant by addictive. The behaviour continues after
the rewards stop, because no number of unrewarded attempts is evidence that the next one will
fail.

It composes exactly with Section~\ref{sec:reward}, and the two accounts are the same account. A
fully predicted reward produces no prediction error. A variable ratio schedule guarantees that no
reward is ever fully predicted. So the schedule is the machinery for keeping the error signal
non-zero indefinitely, and the ramp that peaks at one half is the anticipation it produces while
the player waits to find out.

\subsection{The near miss}

The second finding is sharper and more dangerous. A losing outcome which resembles a win, being
two matching symbols out of three, increases the desire to continue. The effect appears in
experienced gamblers and in people with no gambling history alike, and a near-miss rate of
around thirty per cent is where the effect is reported to be
strongest\footnote{\emph{The Near-Miss Effect in Slot Machines: A Review and Experimental Analysis
Over Half a Century Later}, Journal of the Experimental Analysis of Behavior, 2020.}. Imaging work
finds that near misses recruit reward-related regions in a pattern resembling actual wins.

So thirty per cent is a number the tuner can hold, and near-miss rate is a parameter rather than
an accident.

\subsection{Where the line is}

This material is the mechanism behind loot boxes and behind slot machines, and it is why
Section~\ref{sec:notlicensed} exists. The line is not about whether variance is present. It is
about what the variance is attached to.

\begin{center}
\begin{tabular}{@{}l@{\hspace{5mm}}>{\raggedright\arraybackslash}p{0.58\textwidth}@{}}
\toprule
\textbf{a game}         & The variance resolves a decision. The player chose to go a fourth level
down, and what they find there is the answer to that choice. Playing better shifts the
distribution. \\[1mm]
\textbf{a slot machine} & The variance resolves a repetition. The only input is \emph{again},
skill does not move the distribution, and the near miss is manufactured by stopping the reels
just short. \\
\bottomrule
\end{tabular}
\end{center}

\begin{sectionbox}
Which gives one rule and it is checkable. If the only input to the variable outcome is the
decision to try again, it is a slot machine whatever else is drawn on it. If the distribution the
player is drawing from was selected by their earlier decisions, it is a game, and the intermittent
schedule is doing legitimate work.
\end{sectionbox}

The same test settles the near miss. An honest near miss is informative, meaning the player can
see what they would have needed and can go and get it. A keeper who refuses and names the piece
that was missing is a near miss which teaches. A fabricated near miss conceals that the outcome
was never close, and it is the version that does harm.

\subsection{What this gives the Labyrinth}

Three places where the schedule is already available and one of them is free.

\textbf{Loot.} The depth tables of the design fix what \emph{kind} of thing can be found at each
depth, which is a type-level bound and stays. Which item arrives inside that bound becomes a
draw, and the variance of the draw rises with depth. So push-your-luck and variable ratio become
the same knob, and the deeper level is both more dangerous and more variable.

\textbf{The keeper.} The verdict is already a variable ratio reward and nothing had to be added.
Nothing in the program knows the answer, so the schedule is not simulated, it is the actual
epistemic situation. This is the one reward in the game whose unpredictability is honest all the
way down.

\textbf{The informative refusal.} A keeper who declines and says what was missing is a
thirty-per-cent near miss which also plants an objective. It sends the party back down for a
named thing, which is the loop the design already wanted and had no trigger for.

\begin{alarmbox}
And one thing it must not give. There is no draw anywhere in this game whose only input is a
decision to repeat. Every variable outcome sits behind a choice that cost light, and light is
finite, so the schedule can never become a lever the player pulls indefinitely. That is not a
policy about restraint. It is a structural consequence of having exactly one clock.
\end{alarmbox}

\section{What the genre already knows}

Practice is ahead of theory here, and three games are worth reading closely because each has
already solved part of the problem and one of them documents its own failure.

\subsection{Push your luck, and why intuition fails at it}

The structure is old and the reason it works is arithmetical. Each additional step multiplies the
reward and multiplies the risk, and human estimation of compounded probability is
poor\footnote{Analysis of the structure in Balatro and modern roguelikes: Game Wisdom, \emph{The
Mathematics of Push Your Luck}, 2024.}. So the player is repeatedly asked a question they
cannot answer exactly and can answer approximately, which is the ideal condition for a decision
to feel like judgement rather than calculation.

The Labyrinth already had the shape of this in the phrase about going down as far as you dare,
and had nothing implementing it. Depth was free. Nothing accumulated on the way down and nothing
was risked by continuing.

\subsection{Darkest Dungeon, and the failure it recorded}

The closest published relative of the light mechanic is Darkest Dungeon's light meter. Light
starts at one hundred, depletes as the party moves, at one unit through explored ground and six
through new, and can be restored with a torch or lowered
deliberately\footnote{Darkest Dungeon Wiki, \emph{Light Meter}.}. Lower light raises enemy damage
and accuracy and stress, and raises the quality of loot and the chance of a critical hit.

Two things about it matter for us and they point in opposite directions.

The design succeeds because light is one number that couples to everything. It is described by
players as the closest thing the game has to a difficulty slider, and the decision it poses is
whether the extra profit is worth the extra risk, which is a real question asked continuously.

The design also failed in a specific and instructive way. The intermediate light levels were not
worth occupying, because only the lowest tier paid enough to justify its risk. So the choice
collapsed from a continuum to two options, full light or full dark, and everything between them
was a dead state.

\begin{alarmbox}
That is the failure our bindingness criterion in Section~\ref{sec:new} exists to detect. A
mechanic can be fully implemented, fully documented, and occupy no part of anybody's play. It is
a worse condition than being badly balanced, because being badly balanced is visible.
\end{alarmbox}

One difference is worth naming. Darkest Dungeon's light is a voluntary difficulty slider, where
the player chooses to be in the dark for profit. The Labyrinth's light is a clock, where the
count is rounds remaining and spending it is not optional. The clock is the stronger structure
for our purpose, because it makes every round anywhere in the game cost the same thing, which is
what makes the products in the algebra argue with each other.

\subsection{Into the Breach, and perfect information}

Subset Games built Into the Breach in explicit reaction to the randomness of their previous
game, wanting one in which every death felt like the player's own fault, and reached that by
telegraphing every enemy attack\footnote{Subset Games, IGF interview, Game Developer, 2018.}.

The player is told the enemy's move, the damage, the turn order, the whole map and the state of
every resource. What is withheld is only what cannot be known, being the next wave's spawn points
and what the shop will stock. The reported effect is that the player spends less time working out
how the game works and more time working out how to win, and that each turn becomes a small
puzzle with a discoverable best answer.

\begin{sectionbox}
Which is the same point the prediction error account makes, arrived at from the other side. A
player who cannot predict cannot be surprised, and a player who cannot see the arithmetic cannot
predict. Perfect information is not the opposite of tension. It is the precondition for tension,
because it is what allows the player to know they are in trouble.
\end{sectionbox}

\subsection{The tension, and how it resolves}

There is now an apparent contradiction at the centre of the design and it needs settling rather
than skirting.

The position. Skill depth requires that better play produce better outcomes, and better play
requires the player be able to reason about consequences, so the game should be as legible as
Into the Breach.

The objection. The whole point of the Labyrinth is that its monsters are language models which
may lie, may refuse, and need not answer the same threat the same way twice. That is the
opposite of a telegraphed attack, and it is not a defect to be minimised, since the design
document argues it is the reason for building the thing at all.

The resolution is that these are two different kinds of uncertainty and only one of them is
hostile to skill.

\begin{center}
\begin{tabular}{@{}l@{\hspace{5mm}}>{\raggedright\arraybackslash}p{0.60\textwidth}@{}}
\toprule
\textbf{the economy} & light remaining, arrows, hands, distances, what a weapon takes off, what
loot is worth. Exact, visible, and never randomised. This layer is Into the Breach. \\[1mm]
\textbf{the conversation} & what a monster wants, whether it is telling the truth, whether a
keeper is persuaded. Unpredictable by construction, and unpredictable to the program as much as
to the player. \\
\bottomrule
\end{tabular}
\end{center}

So the non-determinism is confined to the layer where it is the product, and the layer that
carries the arithmetic is made perfectly legible precisely \emph{because} the other layer is
not. A player who cannot be sure whether the goblin is lying must at least be certain how many
rounds of light that conversation cost.

\begin{sectionbox}
And the practical difference is a rule about the interface. Every number in the party record is
displayed exactly and continuously, including the count of rounds and the arithmetic of the way
home. None of it is hidden for atmosphere. The atmosphere is in the monsters and it is the only
place it is allowed to be.
\end{sectionbox}

\section{Plural ends}
\label{sec:plural}

Everything above assumed the player wants one thing and the tuner maximises one number. Both
assumptions are wrong, they are wrong in the same way, and correcting them is the largest single
improvement available to either the game or its tuning.

\subsection{A scalar game is a solved game}

A single-objective push-your-luck game has an optimal policy. Not an unknown one, an optimal one,
and once a player has found it the remaining play is arithmetic. Come home with the largest haul
is a scalar objective, so the Labyrinth as described has a best line, and skill depth measures
only the distance a player still is from finding it.

Compare what the accompanying game \texttt{stalin} does. It asks the player to industrialise a
country, and it holds three ends apart from each other. Do not starve the peasants. Win the war
in 1941. Raise agricultural output above what it would have been. All three are paid for out of
one grain harvest, so they compete, and none of them is expressible in the units of the others.

\begin{sectionbox}
Its own documentation states the consequence plainly. Ranked by soldiers dead the best plan is
brutalist mechanisation, and ranked by everyone dead it comes second, because the fifteen it
starved outnumber the one soldier it saved. The game holds both numbers and does not say which is
the real one.
\end{sectionbox}

That is where the strategy is, and scarcity alone does not produce it. Scarcity plus
\emph{incommensurability} produces it. One clock and one goal is an optimisation problem. One
clock and three goals with no exchange rate between them is a position the player has to take a
view on, and taking a view is what playing is.

\subsection{The autonomy argument, completed}

Section~2.4 recorded that perceived autonomy predicts enjoyment, and the algebraic reading given
there was that a choice must have had genuinely available alternatives. Plural ends give the
stronger version.

Under a single objective, every choice has a correct answer, and the player's autonomy is the
freedom to be wrong. Under plural ends there is no correct answer, only a rate at which the
player is willing to trade one end for another, and the choice is therefore theirs in a sense
that no amount of branching provides. A game with one objective and a thousand paths offers less
autonomy than a game with three objectives and ten.

\subsection{The Labyrinth's three ends}

They exist in the design already and are not distinguished, which is why it reads as one
objective with decoration.

\begin{center}
\begin{tabular}{@{}>{\raggedright\arraybackslash}p{0.20\textwidth}@{\hspace{4mm}}>{\raggedright\arraybackslash}p{0.42\textwidth}@{\hspace{4mm}}>{\raggedright\arraybackslash}p{0.24\textwidth}@{}}
\toprule
\textbf{end} & \textbf{what it is} & \textbf{paid for in} \\
\midrule
\textbf{the haul} & Gold and relics carried to the surface and sold. The only end that is
spendable. &
depth, and hands \\[1mm]
\textbf{the labyrinth opened} & Locks opened from below, puzzles solved, wings entered. Permanent,
compounding, and the only end that survives the party's death. &
hands, light, and pieces that cannot be sold \\[1mm]
\textbf{the bottom, before the end} & Reaching depth eight within a fixed number of descents.
A threshold rather than a quantity. &
parties \\
\bottomrule
\end{tabular}
\end{center}

The first two conflict over a mechanism the design already has and does not use. A puzzle piece
cannot be sold and it occupies a hand. So every piece carried up is a unit of haul declined, and
the choice at the bottom of a descent is not inventory management, it is a decision about which
of two ends this trip was for.

The second and third conflict over light. Rounds spent gathering pieces and arguing with a keeper
are rounds not spent descending, and the calendar does not stop while a puzzle is solved.

\begin{sectionbox}
And the third end is the one the design is missing outright. There is no deadline anywhere in the
document, so nothing punishes patience, and an objective nothing punishes is not an objective. A
fixed number of descents is all it takes, and the multi-party section of the design already has
the parties to spend against it.
\end{sectionbox}

\subsection{The trap that makes it work}

What makes \texttt{stalin} more than three sliders is that the costs are coupled in a direction
the player does not expect. Its own summary puts it as a peasant starved in 1932 being a soldier
missing in 1941, so extracting harder to build steel shrinks the army the steel was meant to
protect. Push far enough and you lose the war holding the largest stockpile in Europe.

The Labyrinth has the same trap available and it is exact. A party lost is exploration absent
later. Drive parties hard to open the labyrinth faster and you arrive at the last descent with
nobody left to make it, holding the most thoroughly unlocked labyrinth nobody will ever reach the
bottom of.

\begin{alarmbox}
That coupling is the difference between plural objectives and a menu. Three independent goals with
independent costs is a choice of which slider to raise, and a player picks one and stops thinking.
Three goals whose costs feed into each other with a delay is a position, and a position has to be
reasoned about every turn.
\end{alarmbox}

\subsection{Front breadth, and the result this report turns on}

Plural ends change what the tuner should measure, and they supply the criterion that was missing.

Run the ladder of Section~\ref{sec:depth} and record each run's outcome not as a score but as a
vector, being haul, labyrinth opened, and whether the bottom was reached. Take the non-dominated
set, meaning the runs which no other run beat on every component at once. That set is the game's
Pareto front, and its width is a measurement.

\begin{center}
\begin{tabular}{@{}l@{\hspace{5mm}}l@{\hspace{5mm}}>{\raggedright\arraybackslash}p{0.36\textwidth}@{}}
\toprule
\textbf{depth} & \textbf{breadth} & \textbf{what you have built} \\
\midrule
low  & low  & a coin flip \\
high & low  & a puzzle, with one answer and a difficulty \\
low  & high & a set of arbitrary preferences, none of them earned \\
high & high & a game \\
\bottomrule
\end{tabular}
\end{center}

Skill depth says bad play is punished. Front breadth says that among good play there is no best.
Both are required and neither implies the other, and this is the central claim of this report.
The first draft of the Labyrinth failed on both, and the economy of the design's
Section~4 repairs only the first.

\begin{sectionbox}
The table of eight strategies in \texttt{stalin}'s README is a Pareto front, printed. Five of the
eight win the war, and the ranking by soldiers dead disagrees with the ranking by everyone dead.
That disagreement is not a defect in the scoring. It is the evidence that the front has width,
and it is the thing a scalar score would have concealed.
\end{sectionbox}

\subsection{So the tuner must not scalarise}

The consequence for optimisation is immediate and it contradicts what the existing tooling in
this repository does.

The current \texttt{optimise.py} collapses eight criteria into a lexicographic scalar, being
criteria met first and then tie-breaks, and searches by random restart and coordinate descent. That
was the right thing to build first and it is the wrong thing to keep, because a scalarisation
selects one point on the front and then cannot report that a front existed. The designer is handed
a winner in place of a choice.

The tools for the multi-objective version are standard. Non-dominated sorting with crowding
distance, being NSGA-II, approximates a front directly and maintains spread along
it\footnote{Deb, Pratap, Agarwal and Meyarivan, \emph{A Fast and Elitist Multiobjective Genetic
Algorithm: NSGA-II}, IEEE Transactions on Evolutionary Computation 6(2), 2002.}. The hypervolume
indicator scores a whole front against a reference point, which is what to use when two fronts
must be compared rather than two settings. Multi-objective Bayesian optimisation exists and keeps
the sample efficiency that Section~\ref{sec:cost} says is mandatory here.

\begin{sectionbox}
And there is a reason to prefer an archive over a front, which is that a designer wants to
\emph{see} the tradeoff rather than be given its resolution. Keeping the best setting found in
each region of behaviour space, rather than the best setting overall, is the quality-diversity
approach, and its output is a map of what kinds of game the parameters can produce. That is a more
useful artefact than a single balanced configuration, because choosing among kinds of game is a
judgement nobody should want automated.
\end{sectionbox}

\subsection{What \texttt{stalin} already proved}

Two findings are in this repository, they were paid for in somebody's time, and both support the
recommendation.

The first is in a comment at the top of \texttt{calibrate.ts}. Tuning against a single number was
what produced a game where brutality was strictly dominated. That is the failure mode of
scalarisation, observed, in this codebase, on this kind of game.

The second is subtler and is the stronger argument. Of the eight things \texttt{stalin} is meant
to say at once, seven hold and the eighth does not, and the accompanying analysis concludes that
the reason is a dependency length rather than a number. In Pareto terms, that target lies outside
the achievable region, so no assignment of the twelve constants will reach it and only a
structural change will.

\begin{alarmbox}
A scalar score cannot tell those two situations apart. Seven out of eight looks like a tuning
problem indefinitely, and the search keeps running. A front tells you which targets are
unreachable in principle, because they sit beyond it, and that is a verdict on the design rather
than on the constants. The eighth criterion was diagnosed by hand here. It should have been
readable off a front.
\end{alarmbox}

\section{Interestingness as a measurable quantity}
\label{sec:browne}

The design question having been answered as far as it can be, the engineering question is whether
any of it reduces to a number. It does, and the precedent is stronger than one might expect.

\subsection{The existence proof}

Cameron Browne and Frederic Maire's \emph{Ludi} system generated combinatorial board games as
rule trees, played them against themselves, and scored the results against a battery of
computable aesthetic criteria\footnote{Browne, \emph{Evolutionary Game Design}, SpringerBriefs in
Computer Science, 2011. Browne and Maire, \emph{Automatic Generation and Evaluation of
Recombination Games}, 2010.}. Fifty-seven criteria were implemented. The system evolved a game
called Yavalath, which was subsequently published commercially and is well regarded by human
players.

That result settles the question of feasibility. A machine measuring self-play statistics against
computable criteria produced a game good enough to sell, without a human designer supplying
taste at any point in the loop. What remains is the translation problem.

\subsection{Which criteria survive one player}

Most of Browne's criteria assume two players with a score between them, and a majority of those
do not translate. Drama in his sense is the degree to which the eventual winner suffered a
negative lead, lead change is counted per turn against an opponent, and completion is about wins
against draws. A single-player push-your-luck game has none of that furniture.

Four translate cleanly.

\begin{center}
\begin{tabular}{@{}l@{\hspace{4mm}}>{\raggedright\arraybackslash}p{0.63\textwidth}@{}}
\toprule
\textbf{skill depth}  & Browne's own version is the correlation between play quality and outcome.
The stronger form is Nielsen's and it is treated separately below. \\[1mm]
\textbf{drama}        & Recast as recovery from a bad position against the game rather than an
opponent. The fraction of successful runs which passed through \texttt{dying}, or through less
than half their light. \\[1mm]
\textbf{duration}     & Kept verbatim. Browne penalises both ends, since a game may be tiresome
or trivial, and both are measured in rounds. \\[1mm]
\textbf{uncertainty}  & Kept, as the spread of outcomes at a fixed setting, and connected to the
one-half target of Section~2.2. \\
\bottomrule
\end{tabular}
\end{center}

\subsection{Skill depth as the primary objective}
\label{sec:depth}

The measure to maximise is Nielsen and Togelius's relative algorithm performance
profile\footnote{Nielsen, Barros, Togelius and Nelson, \emph{General Video Game Evaluation Using
Relative Algorithm Performance Profiles}, EvoApplications 2015. See also Nielsen et al.,
\emph{Evolving Game Skill-Depth using General Video Game AI Agents}, CEC 2017.}. The construction
is to run several agents of known and unequal strength on the same game and examine the ordering
and the size of the gaps. Their reported result is that well-designed games show, on average, a
larger performance difference between stronger and weaker algorithms than badly designed ones
do. A related recent construction for tabletop games builds a skill ladder from the win rate of
an agent given twice the thinking budget of another, across a doubling
sequence\footnote{Goodman et al., \emph{Skill Depth in Tabletop Board Games}, 2024.}.

For the Labyrinth the ladder is four policies, and it is a tensor.

\[ \text{ladder} \;=\; \text{random} \;\boxtimes\; \text{greedy} \;\boxtimes\; \text{timid} \;\boxtimes\; \text{thinker} \]

Random types a legal word. Greedy clears every room and descends whenever it can. Timid clears
one room and climbs out. Thinker is Claude playing to win with the previous run logs available.
The operator is a tensor because all four are prompted with the same setting and none sees what
another did, and that independence is what makes the comparison mean anything.

\begin{sectionbox}
The decisive property of this objective is that it cannot be maximised by making the game hard.
Reduce the light far enough and Thinker scores what Random scores, because nothing survives.
Raise it far enough and they agree again, because nothing is at stake. The gap has an interior
maximum, and difficulty measures have no such thing, which is why difficulty is the wrong
objective and this is the right one.
\end{sectionbox}

\subsection{Two criteria the literature does not supply}
\label{sec:new}

The published criteria were written for games whose rules were fixed and whose parameters were
being searched. The Labyrinth has a specific pathology they do not detect, which is the one
Darkest Dungeon's intermediate light levels exhibited.

\textbf{Bindingness.} For each scarce resource, the fraction of runs in which that resource was
the one that ran out first. If light binds in ninety-five runs out of a hundred and hands never
bind at all, then hands are not a mechanic. The rule exists in the document and in the code and
has no effect on anybody's play. A good setting is one in which all four resources bind
sometimes, and this is the numerical form of the design claim that a product is only a decision
when its branches compete for something scarce.

\textbf{Counterfactual regret.} At each sum in the algebra, the difference in outcome between the
side taken and the side declined. Two ways down with the same outcome behind them are a coin. One
that is always better is a trap. The target is a distribution with both signs in it and a mean
away from zero, so that the choice mattered and neither answer was the answer.

\begin{sectionbox}
Regret is normally unmeasurable, because a path cannot be un-taken. It is cheap here for a reason
peculiar to the architecture. The party record is the entire state that crosses a boundary, and a
descent is a prompt sent to a program started fresh on a port, so the same record can be sent
down both ways against two fresh levels and the two responses compared. The tuner can ask what
would have happened and no player can.
\end{sectionbox}

A third candidate follows from Section~2.2 and is offered more tentatively. \textbf{Pivot
probability}, being the estimated success probability of the marginal decision, tuned towards one
half. It is the most direct expression of the uncertainty finding and the least established of
the six.

\section{Optimisation}

Given six numbers to maximise and roughly ten parameters to vary, the question is which optimiser,
and the answer is dictated entirely by what an evaluation costs.

\subsection{Sample cost is the whole problem}
\label{sec:cost}

Ludi ran self-play games in the thousands, which is affordable when a playthrough is a tree
search over a small board. A playthrough of the Labyrinth with a language model behind every
monster is not in that regime by several orders of magnitude.

The nearest well-documented budget comes from Zook, Fruchter and Riedl, who formalise playtesting
as active learning and tune parameters of a shooter against explicit
goals\footnote{Zook, Fruchter and Riedl, \emph{Automatic Playtesting for Game Parameter Tuning
via Active Learning}, arXiv:1908.01417.}. Their regression goal is stated concretely, being that
the player is hit exactly six times per wave, with the objective the squared difference from that
target. Their reported sample requirement is around seventy evaluations for the regression task
and one to two hundred for the classification task. Notably their evaluations were human
playtests rather than simulated agents, which is what forced the sample efficiency.

\begin{sectionbox}
So the working assumption is hundreds of settings and not thousands, with a batch of playthroughs
per setting. That is precisely the regime Bayesian optimisation exists for, and it rules out
anything that needs a large population per generation.
\end{sectionbox}

\subsection{Two kinds of parameter}

The parameters divide, and they divide by which optimiser can touch them.

\begin{center}
\begin{tabular}{@{}l@{\hspace{4mm}}>{\raggedright\arraybackslash}p{0.60\textwidth}@{}}
\toprule
\textbf{numeric} & rounds per torch, deepest level bearing a torch, arrows per quiver, hands,
sword and bow damage, monster closing speed, rooms per level, monsters per room, labyrinth depth,
loot value by depth \\[1mm]
\textbf{prompt}  & what a monster is told it wants, how readily it is told to lie, how strict a
standard a keeper is told to hold, how much a monster is told it knows about the level beneath \\
\bottomrule
\end{tabular}
\end{center}

The first column is a box in $\mathbb{R}^{10}$ or thereabouts, mostly integer-valued, and it
wants Bayesian optimisation with discrete projection. The second column is not in any
$\mathbb{R}^n$ and no covariance matrix will help. A prompt is improved by reading the runs it
produced and rewriting it, which is something a language model does and a numerical optimiser
cannot.

\subsection{The precedent for coupling them}

That coupling is exactly what \emph{RuleSmith} does. It balances a simplified civilisation game
by running multi-agent language model self-play over textual rulebooks, and wrapping Bayesian
optimisation around a multi-dimensional rule space covering faction mechanics, economy,
production and combat\footnote{\emph{RuleSmith: Multi-Agent LLMs for Automated Game Balancing},
arXiv:2602.06232.}. Its objective is minimising win-rate disparity between factions. Two features
of it transfer directly.

The first is adaptive sampling. Promising candidate settings receive more evaluation games for a
more accurate estimate, and exploratory candidates receive fewer, which is the only way a budget
of hundreds buys any confidence at all.

The second is the reported outcome that the adjustments it proposes are interpretable, being
legible rule changes a human designer can read and apply rather than an opaque vector. That
matters here because the Labyrinth's parameters are all things a person can hold an opinion
about, and a tuner whose output cannot be argued with is less useful than one whose output can.

\subsection{Paying for it, and why the container claim earns its keep}

The remaining obstacle is that thousands of playthroughs with Claude in every monster is not
affordable, and the resolution is a property the architecture already claims.

A monster is an interface. It is a finite set of prompts with a response type for each, and the
response type is one of six actions. The party cannot tell from above what is behind that
interface. So put a table behind it. A monster that chooses among the same six actions by rule
presents the same interface, and the rest of the game cannot distinguish the two.

\begin{center}
\begin{tabular}{@{}l@{\hspace{5mm}}l@{\hspace{5mm}}>{\raggedright\arraybackslash}p{0.36\textwidth}@{}}
\toprule
\textbf{tier} & \textbf{behind the interface} & \textbf{what it tunes} \\
\midrule
cheap & a table & every number, thousands of runs, the full ladder \\
dear  & Claude  & every prompt, tens of runs, Thinker only \\
\bottomrule
\end{tabular}
\end{center}

The cheap tier locates the region of the parameter box. The dear tier checks that the region
survives contact with a monster capable of lying, and rewrites whichever prompts caused it not
to. Between the two, sample adaptively.

\begin{sectionbox}
This is the strongest practical argument the container decomposition makes anywhere. The claim
that a language model behind a socket is interchangeable with a table behind a socket is not a
piece of architectural tidiness. It is the difference between a tuning run that costs a few
dollars and one that cannot be run.
\end{sectionbox}

\newpage

\section{Recommendations, with confidence}

Ranked by expected effect on the two complaints, with an honest estimate against each. The design
changes come first because the tuning changes are worthless without them.

\small
\begin{center}
\begin{tabular}{@{}>{\raggedright\arraybackslash}p{0.24\textwidth}@{\hspace{3mm}}>{\raggedright\arraybackslash}p{0.50\textwidth}@{\hspace{3mm}}>{\raggedright\arraybackslash}p{0.15\textwidth}@{}}
\toprule
\textbf{change} & \textbf{why} & \textbf{confidence} \\
\midrule
Three ends held apart &
A scalar game has an optimal policy and is therefore solved. Scarcity produces arithmetic, and
scarcity plus incommensurability produces a position. \texttt{stalin} is the working precedent. &
high \\[2mm]
A deadline, in descents &
The third end does not exist in the design at all, and an objective nothing punishes is not an
objective. &
high \\[2mm]
Pareto front, never a scalar score &
\texttt{calibrate.ts} records that tuning against a single number is what produced a game where
brutality was strictly dominated. Observed, in this codebase. &
high \\[2mm]
Front breadth as a second objective &
Skill depth alone selects for puzzles. Breadth is what distinguishes a game from a solved one.
Novel, and untested. &
medium-high \\[2mm]
Variable ratio on the loot draw, inside the depth-typed bound &
The most extinction-resistant schedule there is, and the type-level bound keeps it honest. The
keeper's verdict is already one and cost nothing. &
high \\[2mm]
Informative refusal at around thirty per cent &
The near-miss rate the literature reports as strongest, and the informative form plants an
objective instead of concealing one. &
medium \\[2mm]
One scarce clock, being light &
Supplies the prediction that reward error needs, makes every product compete, and turns depth
into arithmetic. Darkest Dungeon is the precedent and its failure mode is known. &
high \\[2mm]
Rooms as a product, not a tensor &
An independent conjunction is a chore with an operator on it. This is the autonomy requirement
in algebraic form. &
high \\[2mm]
Perfect information on the economy &
Into the Breach's result and the prediction error account agree. Legibility is the precondition
for tension rather than its enemy. &
high \\[2mm]
\bottomrule
\end{tabular}
\end{center}

\begin{center}
\begin{tabular}{@{}>{\raggedright\arraybackslash}p{0.24\textwidth}@{\hspace{3mm}}>{\raggedright\arraybackslash}p{0.50\textwidth}@{\hspace{3mm}}>{\raggedright\arraybackslash}p{0.15\textwidth}@{}}
\toprule
\textbf{change} & \textbf{why} & \textbf{confidence} \\
\midrule
Skill depth as the tuning objective &
Published, validated, and cannot be gamed by raising difficulty. &
high \\[2mm]
Bindingness as a guard &
Detects the exact failure Darkest Dungeon recorded. Cheap to compute. Not in the literature, so
untested. &
medium-high \\[2mm]
Death costing a party and not a labyrinth &
Standard in the genre and free here, since levels were being kept anyway. &
medium-high \\[2mm]
Two-tier tuning behind a fake monster &
Rests on the container claim, which is the document's central thesis. If it fails, the thesis
was wrong about something more important than tuning. &
medium-high \\[2mm]
Monsters knowing one true fact &
Converts the way down from a coin flip into a decision and makes the non-determinism
load-bearing. Untested and the most likely to need iteration. &
medium \\[2mm]
Counterfactual regret &
Uniquely cheap here and genuinely novel. No precedent, so no idea whether the distribution it
produces is informative. &
medium \\[2mm]
Pivot probability near one half &
Best-supported psychology and the longest inferential chain. Adopt as a tiebreak, not a primary
objective. &
low-medium \\
\bottomrule
\end{tabular}
\end{center}
\normalsize

\section{What would falsify this}

Three things would, and each is checkable early.

If the ladder does not separate at any setting, meaning Thinker never reliably beats Greedy
anywhere in the parameter box, then the game has no skill in it and no amount of tuning will
install any. That is a verdict on the mechanics rather than on the numbers, and it is the first
thing the cheap tier should be asked.

If bindingness cannot be brought away from a single resource, meaning light binds almost always
regardless of the setting, then three of the four scarce things are decoration and the honest
response is to remove them rather than to keep tuning.

If the region found by the cheap tier does not survive the dear tier, meaning the game balanced
against table-driven monsters plays entirely differently against monsters that can lie, then the
container claim has failed where it matters, and the tuning must be done at a cost that may not
be payable.

And one more, added with Section~\ref{sec:plural} and arguably the most important of the four. If
the front collapses to a point at every setting, meaning one way of playing dominates on all three
ends everywhere in the parameter box, then the three ends are not really three. Either they do not
genuinely compete, or one of them is a proxy for another, and the fix is in the design rather than
in the constants. This is the failure \texttt{stalin}'s eighth criterion turned out to be, and it
was diagnosed by hand there because there was no front to read it off.

\begin{sectionbox}
The first of the three is the one to run first, because it is the cheapest and it is the only one
whose failure would mean the design is wrong rather than the numbers.
\end{sectionbox}

\section*{Sources}

\small
\begin{itemize}\setlength{\itemsep}{1pt}
\item Fiorillo, Tobler and Schultz. Discrete Coding of Reward Probability and Uncertainty by
Dopamine Neurons. \emph{Science} 299, 2003.
\item Schultz. Dopamine reward prediction-error signalling: a two-component response.
\emph{Nature Reviews Neuroscience} 17, 2016.
\item Ryan, Rigby and Przybylski. The Motivational Pull of Video Games. \emph{Motivation and
Emotion} 30, 2006.
\item Przybylski, Rigby and Ryan. A Motivational Model of Video Game Engagement. \emph{Review of
General Psychology} 14, 2010.
\item Dah et al. Gamification is not Working: Why? \emph{Games and Culture}, 2025.
\item Browne. \emph{Evolutionary Game Design}. Springer, 2011.
\item Browne and Maire. Automatic Generation and Evaluation of Recombination Games, 2010.
\item Nielsen, Barros, Togelius and Nelson. General Video Game Evaluation Using Relative
Algorithm Performance Profiles. \emph{EvoApplications}, 2015.
\item Nielsen et al. Evolving Game Skill-Depth using General Video Game AI Agents. \emph{CEC},
2017.
\item Goodman et al. Skill Depth in Tabletop Board Games, 2024.
\item Zook, Fruchter and Riedl. Automatic Playtesting for Game Parameter Tuning via Active
Learning. arXiv:1908.01417.
\item RuleSmith: Multi-Agent LLMs for Automated Game Balancing. arXiv:2602.06232.
\item Deb, Pratap, Agarwal and Meyarivan. A Fast and Elitist Multiobjective Genetic Algorithm:
NSGA-II. \emph{IEEE Transactions on Evolutionary Computation} 6(2), 2002.
\item Ferster and Skinner. \emph{Schedules of Reinforcement}, 1957.
\item The Near-Miss Effect in Slot Machines: A Review and Experimental Analysis Over Half a
Century Later. \emph{JEAB}, 2020.
\item This repository. \texttt{calibrate.ts}, \texttt{optimise.py}, \texttt{README.md}, and
\texttt{implementation.pdf}.
\item Subset Games. Road to the IGF interview. \emph{Game Developer}, 2018.
\item Darkest Dungeon Wiki. Light Meter.
\item Game Wisdom. The Mathematics of Push Your Luck, 2024.
\end{itemize}
\normalsize

\end{document}
